% Authored. Tables and numbers come from generated/.

\section{Robustness and Accounting Boundaries}
\label{sec:robustness}

\subsection{The 1995 splice}

The splice is the single largest threat to the level results, and it is handled by
restriction rather than by adjustment. Both 1995 vintages are retained; the revision
between them is non-zero for every subsector; no level statistic is averaged across
the boundary; and no annual change is computed through it. The regime split of
Table~\ref{tab:composition} is the form in which magnitudes are reported precisely
because a pooled magnitude would embed the splice silently.

Sign-based results are less exposed to the splice than magnitudes, and
Section~\ref{sec:data} checks that rather than asserting it: sign classifications are
unchanged across both source overlaps this panel retains. The persistence results of
Section~\ref{sec:results} therefore carry more weight than the means beside them, and
we have kept the two visually separate for that reason.

\subsection{How firm are the change-point dates?}

The piecewise-constant mean model of Section~\ref{sec:method} selects break dates,
and a selected date reads as more determined than it is --- particularly here, where
the penalised criterion follows \citet{yao1988} in form but is one of several the
change-point literature admits \citep{baiperron1998}. Two guards are reported.

The first is a BIC margin over the next-best break count, which says how decisively
the selection beat its alternatives. The second is a sensitivity grid: detection
re-run over \BreakSpecifications\ combinations of the two tuning parameters, neither
of which is estimated from the data. A date that survives only one combination is a
property of that choice rather than of the series.

Table~\ref{tab:breakstability} reports the grid summary, and
Appendix~\ref{sec:changepoints} the full detail.

\begin{table}[htbp]
\centering
\caption{Sensitivity of the detected mean shifts to the two tuning choices, over twelve specifications per series.}
\label{tab:breakstability}
\small
\fitwidth{%
\begin{tabular}{@{}llrlrr@{}}
\toprule
Regime & Subsector & Modal breaks & Modal dates & Share at modal dates & Distinct date sets \\
\midrule
1977--1994 historical & General Government & 1 & 1986 & 0.83 & 2 \\
1995--2025 modern & General Government & 2 & 2009, 2015 & 0.42 & 4 \\
1977--1994 historical & Central Government & 1 & 1987 & 1.00 & 1 \\
1995--2025 modern & Central Government & 2 & 2009, 2016 & 0.67 & 2 \\
1977--1994 historical & Regional and Local & 0 & -- & 0.83 & 2 \\
1995--2025 modern & Regional and Local & 2 & 2012 & 0.33 & 6 \\
1977--1994 historical & Social Security Funds & 0 & -- & 1.00 & 1 \\
1995--2025 modern & Social Security Funds & 1 & 2017 & 0.67 & 3 \\
\bottomrule
\end{tabular}%
}
\end{table}


The grid is not a formality: it changes what may be claimed. In
\BreakModalAgree\ of the \BreakSeries\ series the preferred specification returns the
modal set of dates. In the remaining series it does not, and for those the selected
dates follow the tuning choice. The least stable case is
\WeakestBreakSector\ in the modern regime, where the modal dates hold in only
\WeakestBreakShare\% of the grid across \WeakestBreakSets\ distinct date sets.

Accordingly we describe every date as a candidate rather than a finding, and no
break year should be quoted from this paper without its share. We also attach no
historical cause to any of them, and note that a five-year minimum segment means a
shift in the last four years of the sample cannot be detected at all.

\subsection{Intergovernmental transfers and the location of a balance}

Because the subsectors are consolidated, a transfer between them changes the
components and not the aggregate. This makes the \emph{location} of a balance within
general government partly a function of the transfer convention, which matters for
interpreting composition results.

The historical source tables identify current and capital transfers between public
administrations, which supports a purely mechanical sensitivity for
\PanelStart--1995:
\[
B^{sens}_{i,t} = B_{i,t} - \left(T^{received}_{i,t} - T^{paid}_{i,t}\right).
\]

We compute it and deliberately decline to label it an underlying or true balance,
and we never use it to restate B.9. The reason is substantive rather than cautious:
a transfer normally finances an expenditure responsibility assigned to the
recipient, so removing the transfer while leaving the responsibility in place does
not describe a coherent alternative arrangement. Which transfer discharges which
statutory obligation is a legal question, not an accounting one. What the
sensitivity quantifies is how much the recorded location of a balance depends on the
convention --- and it is reported for that purpose only.

The same reasoning is why no Social Security balance in this paper has State
transfers netted out of it.

\subsection{What the balance does not determine}

Two limits on the reach of these results are worth stating explicitly, because both
concern inferences a reader may reasonably be tempted to draw.

The balance does not determine the debt path. The change in debt equals minus the
balance plus a stock-flow adjustment that absorbs financial-asset transactions,
valuation effects and timing differences, and in several years in these data the
adjustment is the larger of the two terms. Whether such residuals reflect accounting
responses to fiscal rules \citep{vonhagen2006} or are smaller and less systematic
than that literature suggested \citep{seiferling2013} is a live question on which we
take no position. Its relevance here is narrower: a composition result about
balances is not a result about debt.

The balance does not identify a mechanism. The decompositions of
Section~\ref{sec:method} locate a movement in the accounts; they do not explain it.
Nothing here separates cyclical from discretionary movements, and no cyclical
adjustment is attempted. A specification relating balances to nominal GDP growth exists
in the repository and is reported there with its weaknesses set out. It supports no claim
in this paper, because its dependent variable and its regressor share a construction
through nominal GDP and its fits are very low --- \GdpRSquaredLow\ to
\GdpRSquaredHigh\ across \GdpRSquaredSeries\ series. It appears here for one purpose
only, as the point of comparison in Section~\ref{sec:contributionbase}: quoting its fit
to say that a better-specified regressor does better is a use of the number that its
weaknesses do not undermine.

For the Social Security side specifically, that gap can be narrowed rather than left open.
Equation~\eqref{eq:ssfchange} shows that contributions are the dominant positive term in
the recent balance changes, and the Council attributes the contribution growth to rising
employment and average earnings \citep{cfp_ss_outturn2025}.
Section~\ref{sec:contributionbase} models the contribution base directly, relating the
change in contributions to the change in the aggregate wage bill $W_t = N_t \bar{w}_t$
rather than to aggregate nominal growth. The regressor there is the largest observable
reference base for the levy, and at \WageBillRSquared\ the fit is several times tighter
than anything the specification criticised here achieves.
